\chapter{第二十四讲\quad 渎职罪}
\label{lecture:24}
\noindent{\small\hyperref[ref:lecture:24]{查阅本讲扩展规则}}\par

\section{第一节\quad 重点罪名}

\paragraph{法条与归责模型}
设 \(I\) 为合格主体，\(D\) 为具体职责，\(V\) 为违反职责，\(L\) 为重大损失，\(C\) 为该违反与损失之间的因果归责。第397条一般入口核对 \(I\land D\land V\land L\land C\) 及相应罪过。\(V\) 在滥用职权分支是超越职权或不正当行使职权等，在玩忽职守分支是严重不负责任而不履行或不正确履行职责。区分两罪时应分别核对行为和结果的主观态度，不能仅以有动作还是无动作判断。法释〔2012〕18号第2、7条要求另查具体渎职罪优先适用与依法或者受委托行使国家行政管理职权的主体范围。

\paragraph{推导一}
负有依法查封危险设施职责的国家机关工作人员甲，严重不负责任地遗漏依法必须作出的查封。证据确定如按职责处理即可避免该设施随后导致的一人死亡，且没有其他切断归责的原因。甲的身份及法定职责给出 \(I,D\)，严重遗漏给出玩忽职守型 \(V\)，一人死亡依2012年渎职解释第1条给出 \(L\)，死亡实现了查封义务针对的风险，结合事实因果得 \(C\)。若无更具体条款，依第397条审查玩忽职守罪。把死亡改为与设施无关的突发疾病，便不能沿用原有 \(C\)。

\paragraph{推导二}
司法工作人员乙明知被调查者无罪，为徇私使其受刑事追诉。身份、明知、徇私和使其受追诉分别得到证实，依第399条第一款得徇私枉法罪的候选。民事、行政枉法裁判按该条第二款另核故意违背事实和法律作裁判及情节严重。执行判决、裁定中的失职或滥用职权按第三款另核具体措施、职责及重大损失。若乙的第一款行为同时满足第385条受贿罪，优先用第399条第四款。
\[
B_{399,1}\land B_{385}
  \xrightarrow{\text{第399条第四款}}\mathrm{SelectHeavier}(399,385).
\]
\(\mathrm{SelectHeavier}\) 指按处罚较重的规定定罪处罚。

\paragraph{边界}
其他渎职行为同时构成受贿罪时，2012年渎职解释第3条及2016年贪污贿赂解释第17条规定，在各自适用范围内除刑法另有规定外数罪并罚。第399条第四款是特别规则，不能推广到全部渎职罪。

\section{第二节\quad 普通罪名}

\paragraph{法条与职责索引}
将职责 \(D\) 细分为看押、移送、查禁、解救、食药监管五类。第400条故意私放在押人员与严重不负责任致其脱逃并造成严重后果，属于不同分支。第402条需行政执法人员、依法应移交、徇私舞弊而不移交与严重情节。第417条以有查禁犯罪职责者通风报信或提供便利帮助犯罪分子逃避处罚为核心。第416条区分接到解救要求或举报而不解救并造成严重后果，以及利用职务阻碍解救。第408条之一核对食药安全监管职责、法定滥权或失职行为、严重后果或其他严重情节。根据罪名补充规定（七），其正式罪名是「食品、药品监管渎职罪」，其两类行为可以分设谓词。

\paragraph{推导}
司法工作人员甲故意打开牢门，让依法在押的被告人离开监管控制。角色、羁押对象和主动放行事实给出第400条第一款的主体、对象和私放行为，依该款审查。若甲因严重疏忽忘记上锁，嫌疑人乘机脱逃，则「故意私放」谓词失效，改查第二款的严重不负责任、因果关系及严重后果。后果未证明时，第二款的合取式不能完成。

\paragraph{边界}
行政执法人员面对尚不具备依法移交条件的一般违法，不移交不能仅因此推出第402条。刑罚执行中的外部人员只因知晓位置也没有第400条所需的司法工作人员身份，须另查是否与具备身份者存在共同犯罪。

\paragraph{闭卷记忆}
\begin{enumerate}
\item 先记「身、责、违、果、因」，依次写 \(I,D,V,L,C\)，再核该罪的罪过。职责与归责必须各写一句证据。
\item 特殊职责画「押、送、禁、救、食药」，分别定位400、402、417、416、408之一。定位后再查各条情节。
\item 受贿另列一行。先检查第399条第四款等特别规则，再决定择重或并罚。
\end{enumerate}

\paragraph{闭卷自测与答案}
\begin{enumerate}
\item 执法人员没有依法移交，能直接定第402条吗？答：还需法定主体、依法应移交、徇私舞弊及严重情节，不能以结果动作代替其余要件。
\item 徇私枉法同时构成受贿，是否适用一般并罚规则？答：若符合第399条第四款，适用特别择重规则。
\end{enumerate}
